\section{Dense radial singularities}
\label{sec:natural}

Keep the binary golden control and set \(E(0)=0\).  The variable in
\[
  G(z)=\sum_{N\ge0}E(N)z^N
\]
marks the prefix cutoff \(N\).  Because \cref{thm:qadic-boundary} makes
\(E(N)\) bounded, \(G\) is analytic on the open unit disk.

For an integer \(Q\ge2\), periodicity gives the rational residue series
\begin{equation}
  R_Q(z):=\sum_{N\ge0}(N\bmod Q)z^N
  =\frac{P_Q(z)}{1-z^Q},\qquad
  P_Q(z)=\sum_{a=0}^{Q-1}az^a.
  \label{eq:residue-generating}
\end{equation}
Writing \(\Delta_v=d_v-d_{v-1}\), absolute convergence of
\(\sum_v|\Delta_v|\) permits the exact rearrangement
\begin{equation}
  G(z)=-\sum_{v\ge1}\frac{\Delta_v}{2^v}R_{2^v}(z),
  \qquad |z|<1.
  \label{eq:G-residue}
\end{equation}

\begin{theorem}[Radial coefficients and natural boundary]
\label{thm:natural}
Let \(\xi\) be a primitive \(2^v\)-th root of unity, \(v\ge1\).  Then
\begin{equation}
  \lim_{s\uparrow1}(1-s)G(s\xi)
  =-\frac{\gamma_{v-1}}{2^{v-1}(1-\xi)}\ne0.
  \label{eq:radial-coefficient}
\end{equation}
The unit circle is a natural boundary for \(G\).
\end{theorem}

\begin{proof}
If \(w<v\), then \(1-\xi^{2^w}\ne0\), so the radial leading coefficient
of \(R_{2^w}\) vanishes.  If \(w\ge v\), set \(Q=2^w\).  Since
\(\xi^Q=1\) and \(\xi\ne1\), differentiating the finite geometric sum gives
\[
  P_Q(\xi)=-\frac{Q}{1-\xi}.
\]
Together with \(1-s^Q\sim Q(1-s)\), this implies
\begin{equation}
  \lim_{s\uparrow1}(1-s)R_Q(s\xi)=-\frac1{1-\xi}.
  \label{eq:residue-radial}
\end{equation}
The limit may pass through \eqref{eq:G-residue}.  Indeed,
\[
  \frac{1-s}{Q}|R_Q(s\xi)|
  \le\frac{1-s}{Q}\sum_{N\ge0}Qs^N=1,
\]
so the level-\(w\) contribution is dominated by \(|\Delta_w|\).  Dominated
convergence now yields
\[
  \lim_{s\uparrow1}(1-s)G(s\xi)
  =\frac1{1-\xi}\sum_{w\ge v}\frac{\Delta_w}{2^w}.
\]
By \eqref{eq:gamma-residue},
\(\sum_{w\ge v}\Delta_w/2^w=-\gamma_{v-1}/2^{v-1}\), which proves
\eqref{eq:radial-coefficient}; nonvanishing follows from
\cref{thm:golden}.

Primitive dyadic roots are dense on the unit circle.  If \(G\) admitted
analytic continuation across any boundary point, the continuation domain
would contain an open boundary arc and hence a primitive dyadic root
\(\xi\).  More explicitly, let \(H\) be holomorphic on a neighborhood
\(U\ni\xi\) and agree with \(G\) on \(U\cap\mathbb D\), as forced by the
identity theorem.  Holomorphy makes \(H\) bounded on some closed disk about
\(\xi\) contained in \(U\).  Equation \eqref{eq:radial-coefficient} instead
makes \(G(s\xi)=H(s\xi)\) unbounded like a nonzero multiple of
\((1-s)^{-1}\) as \(s\uparrow1\), a contradiction.
\end{proof}

\paragraph{Normalization control.}
The first nonreal control prevents a silent phase error.  For \(Q=4\) and
\(\xi=i\),
\[
  P_4(i)=-2-2i=-\frac4{1-i},\qquad
  \lim_{s\uparrow1}(1-s)R_4(si)
  =-\frac1{1-i}=\frac{-1-i}{2}.
\]
This differs from \(i/(1-i)=(-1+i)/2\).  The exact coefficient in
\eqref{eq:radial-coefficient}, rather than a merely proportional
expression, is therefore essential.

\begin{figure}[t]
  \centering
  \resizebox{\textwidth}{!}{\begin{tikzpicture}[
  >=Latex,
  every node/.style={font=\small},
  box/.style={rounded corners,draw,thick,align=center,text width=3.35cm,
              minimum height=1.45cm,inner sep=5pt}
]
  \node[box,draw=chainblue] (rem) at (0,1.25)
    {\(\displaystyle E(N)=
      -\sum_{v\ge1}\frac{\Delta_v}{2^v}(N\bmod2^v)\)};
  \node[box,draw=boundaryteal] (rat) at (5.0,1.25)
    {\(\displaystyle R_{2^v}(z)=
      \frac{\sum_{a=0}^{2^v-1}az^a}{1-z^{2^v}}\)};
  \node[box,draw=goldenorange,text width=3.7cm] (rad) at (10.0,1.25)
    {\(\displaystyle \lim_{r\uparrow1}(1-r)G(r\xi)
      =-\frac{\gamma_{v-1}}{2^{v-1}(1-\xi)}\ne0\)};
  \node[box,draw=proofgreen] (nat) at (15.0,1.25)
    {primitive dyadic roots are dense\\
     \(\Longrightarrow\) unit circle is a natural boundary};
  \draw[->,very thick] (rem)--
    node[above,font=\scriptsize,fill=white,inner sep=1pt] {residue OGF}(rat);
  \draw[->,very thick] (rat)--
    node[above,font=\scriptsize,fill=white,inner sep=1pt] {radial limit}(rad);
  \draw[->,very thick] (rad)--
    node[above,font=\scriptsize,fill=white,inner sep=1pt] {dense set}(nat);

  \node[box,draw=stopred,fill=stopred!5,text width=10.2cm,
        minimum height=1.35cm] (forbid) at (7.5,-1.25)
    {\textcolor{stopred}{\textbf{Type firewall.}}
     \(N\) is a prefix cutoff, not a periodic orbit; \(G\) is neither an
     Artin--Mazur zeta nor a Fredholm determinant; dense radial singularities
     are not isolated meromorphic poles.};
\end{tikzpicture}
}
  \caption{The analytic implication and its type firewall.  Residue
  generating functions yield nonzero radial leading coefficients at a dense
  set of dyadic roots, which obstruct continuation across every boundary
  arc.  The marker is a prefix cutoff: no orbit zeta, trace, determinant, or
  isolated-pole interpretation is used.}
  \label{fig:radial-firewall}
\end{figure}

\begin{remark}[Radial is not meromorphic]
Equation \eqref{eq:radial-coefficient} is an Abelian radial statement.
Because the singular boundary points are dense, they are not isolated poles
of a meromorphic continuation.  The conclusion is precisely the natural
boundary and no stronger meromorphic assertion.
\end{remark}
